\bta{2003}


\begin{enumerate}
	%\renewcommand{\labelenumi}{\arabic{enumi}.}
	% A(\Alph) a(\alph) I(\Roman) i(\roman) 1(\arabic)
	%设定全局标号series=example	%引用全局变量resume=example
	%[topsep=-0.3em,parsep=-0.3em,itemsep=-0.3em,partopsep=-0.3em]
	%可使用leftmargin调整列表环境左边的空白长度 [leftmargin=0em]
	\item
	填空题 (本题共 5 小题，每小题 5 分，满分 25 分)
\begin{enumerate}
	%\renewcommand{\labelenumi}{\arabic{enumi}.}
	% A(\Alph) a(\alph) I(\Roman) i(\roman) 1(\arabic)
	%设定全局标号series=example	%引用全局变量resume=example
	%[topsep=-0.3em,parsep=-0.3em,itemsep=-0.3em,partopsep=-0.3em]
	%可使用leftmargin调整列表环境左边的空白长度 [leftmargin=0em]
	\item
	$\lim\limits _{x \rightarrow \frac{\pi}{4}} \frac{\sqrt[3]{\tan x}-1}{2 \sin ^{2} x-1}$ \tk{$  \frac{ 1 }{ 3 }  $} 。
	\item 
	设 $\left\{\begin{aligned}x=\int_{1}^{t^{2}} u \ln u d u \\ y=\int_{t^{2}}^{1} u^{2} \ln u d u\end{aligned}\right.$ $(t>0)$, 则 $\frac{d^{2} y}{d x^{2}}=$ \tk{$ \frac{-1}{2t^{2}\ln t} $} 。
	\item 
	级数 $\sum\limits_{n=0}^{\infty} \frac{1}{2 n+1}\left(\frac{1-2 x}{1+x}\right)^{n}$ 的收敛域为 \tk{$ (0,2] $} 。
	\item 
椭球面 $x^{2}+2 y^{2}+z^{2}=1$ 上平行于平面 $x-y+2 z=0$ 的切平面方程为 $x-y+2 z=$ \tk{$ \pm\sqrt{\frac{11}{2}} $} 。
	
	
	\item
	微分方程 $y^{\prime \prime}- y^{\prime}-2 y=4 x e^{x}$ 的通解为 \tk{$ y=-(2x+1)e^{x}+C_{1}e^{-x}+C_{2}e^{2x} $} 。
	
	
\end{enumerate}

\banswer{
\begin{enumerate}
	%\renewcommand{\labelenumi}{\arabic{enumi}.}
	\item
	$\frac13$。令 $t=\tan x$，则
	\[
	2\sin^{2}x-1=\frac{2t^{2}}{1+t^{2}}-1=\frac{t^{2}-1}{t^{2}+1},
	\]
	故原式 $=\dfrac{\left(\sqrt[3]{t}-1\right)\left(t^{2}+1\right)}{(t-1)(t+1)}$。
	当 $x\to\frac\pi4$ 时 $t\to1$，而
	\[
	\lim_{t\to1}\frac{\sqrt[3]{t}-1}{t-1}=\frac13,\qquad
	\lim_{t\to1}\frac{t^{2}+1}{t+1}=1 ,
	\]
	所以极限为 $\frac13$。

	\item
	$-\dfrac{1}{2t^{2}\ln t}$。由变限积分求导（注意 $\ln t^{2}=2\ln t$）
	\[
	\frac{dx}{dt}=t^{2}\ln t^{2}\cdot2t=4t^{3}\ln t,\qquad
	\frac{dy}{dt}=-\left(t^{2}\right)^{2}\ln t^{2}\cdot2t=-4t^{5}\ln t ,
	\]
	故 $\dfrac{dy}{dx}=-t^{2}$。再对 $x$ 求导一次：
	\[
	\frac{d^{2}y}{dx^{2}}
	=\frac{\frac{d}{dt}\left(-t^{2}\right)}{\frac{dx}{dt}}
	=\frac{-2t}{4t^{3}\ln t}=-\frac{1}{2t^{2}\ln t}\qquad(t>0,\ t\neq1).
	\]

	\item
	$(0,2]$。记 $u=\dfrac{1-2x}{1+x}$。幂级数 $\sum\limits_{n=0}^{\infty}\frac{u^{n}}{2n+1}$
	的收敛半径为 $1$；$u=-1$ 时是收敛的交错级数，$u=1$ 时是发散的
	$\sum\frac{1}{2n+1}$，故收敛条件为 $-1\leqslant u<1$。由
	\[
	u-1=\frac{-3x}{1+x},\qquad u+1=\frac{2-x}{1+x}
	\]
	知 $u<1\iff\dfrac{x}{1+x}>0$，$u\geqslant-1\iff\dfrac{2-x}{1+x}\geqslant0$，
	两者同时成立当且仅当 $0<x\leqslant2$。

	\item
	$x-y+2z=\pm\sqrt{\frac{11}{2}}$。设切点为 $\left(x_{0},y_{0},z_{0}\right)$，
	$F=x^{2}+2y^{2}+z^{2}-1$ 的梯度 $\left(2x_{0},4y_{0},2z_{0}\right)$
	应与已知平面的法向量 $(1,-1,2)$ 平行，故
	\[
	x_{0}=\frac\lambda2,\qquad y_{0}=-\frac\lambda4,\qquad z_{0}=\lambda .
	\]
	代入椭球面方程得 $\frac{\lambda^{2}}{4}+\frac{\lambda^{2}}{8}+\lambda^{2}
	=\frac{11\lambda^{2}}{8}=1$，即 $\lambda=\pm\sqrt{\frac{8}{11}}$。
	切平面为 $x_{0}x+2y_{0}y+z_{0}z=1$，即 $\frac\lambda2(x-y+2z)=1$，于是
	\[
	x-y+2z=\frac{2}{\lambda}=\pm\sqrt{\frac{11}{2}}=\pm\frac{\sqrt{22}}{2}
	\approx\pm2.3452 .
	\]

	\par\vspace{1mm}\noindent
	\textbf{注：}{\kaishu
	原答案作 $x-y+2z=\pm\sqrt{\frac{y}{2}}$，右端竟含变量 $y$，显系 $11$ 的排印之误，
	今改正为 $\pm\sqrt{\frac{11}{2}}$。已数值校验：上述切点确在椭球面上，
	且其梯度与 $(1,-1,2)$ 成比例。另外原题面“的且平面方程”应为
	“的切平面方程”，一并改正。
	}\par

	\item
	$y=-(2x+1)e^{x}+C_{1}e^{-x}+C_{2}e^{2x}$。特征方程 $r^{2}-r-2=0$ 的根为
	$r_{1}=-1$、$r_{2}=2$，故齐次通解为 $C_{1}e^{-x}+C_{2}e^{2x}$。
	因 $1$ 不是特征根，设特解 $y^{*}=(ax+b)e^{x}$，则
	\begin{align*}
	\left(y^{*}\right)^{\prime\prime}-\left(y^{*}\right)^{\prime}-2y^{*}
	&=(-2ax+a-2b)e^{x}=4xe^{x},
	\end{align*}
	比较系数得 $a=-2$、$b=-1$，即 $y^{*}=-(2x+1)e^{x}$。

	\par\vspace{1mm}\noindent
	\textbf{注：}{\kaishu
	原题面印作 $y^{\prime\prime}+2y^{\prime}-2y=4xe^{x}$，与所给答案对不上：
	该方程的特征根为 $-1\pm\sqrt3$，通解应为
	$(4x-16)e^{x}+C_{1}e^{\left(-1+\sqrt3\right)x}+C_{2}e^{\left(-1-\sqrt3\right)x}$，
	齐次部分与特解都不吻合。而答案 $-(2x+1)e^{x}+C_{1}e^{-x}+C_{2}e^{2x}$
	逐项恰好是 $y^{\prime\prime}-y^{\prime}-2y=4xe^{x}$ 的通解（已代入数值验证），
	故判定题面中的 $+2y^{\prime}$ 系 $-y^{\prime}$ 之误，已据此更正题面。
	}\par
\end{enumerate}
}

\item
选择题 (本题 5 小题，每小题 5 分，满分 25 分)
\begin{enumerate}
	%\renewcommand{\labelenumi}{\arabic{enumi}.}
	% A(\Alph) a(\alph) I(\Roman) i(\roman) 1(\arabic)
	%设定全局标号series=example	%引用全局变量resume=example
	%[topsep=-0.3em,parsep=-0.3em,itemsep=-0.3em,partopsep=-0.3em]
	%可使用leftmargin调整列表环境左边的空白长度 [leftmargin=0em]
	\item
	设$f(x)=\left\{\begin{array}{ll}|x|^{a} \arctan \frac{1}{x} & x \neq 0 \\ x & x=0\end{array}\right.$在 $x=0$ 处连续但不可导, 则 $a$的取值范围是 \xzanswer{C} 
	
	
\fourchoices
{$a>0$}
{$0<a \leqslant 1$}
{$0<a<1$}
{$a>1$}

\item 


\item
“对任意给定的 $\epsilon>0$, 总存在正整数 $N$, 使得当 $n>N$ 时，就有 $\mid a_{N+1}+a_{N+2}+\cdots+a_{n} \mid<\epsilon $” 是级数 $\sum\limits_{n=1}^{\infty} a_{n}$ 收敛的 \xzanswer{C} 


\fourchoices
{充分条件但非必要条件}
{必要条件但非充分条件}
{充分必要条件}
{既非充分条件也非必要条件}

	
\item 
设 $S$ 是柱面 $x^{2}+y^{2}=R^{2}$ 介于平面 $z=0$ 及 $z=R$ 之间的部分, 则 $\iint_{S}\left(x^{2}+z^{2}\right) d S=$ \xzanswer{B}


\fourchoices
{$\frac{8}{3} \pi R^{4}$}
{$\frac{5}{3} \pi R^{4}$}
{$\frac{4}{3} \pi R^{4}$}
{$\pi R^{4}$}

\item 
设 $L$ 是起点 $A(-1,0)$, 终点 $B(1,0)$ 的简单光滑曲线, 除 $A, B$ 外其他点都在 $x$ 轴
上方, 则曲线积分 $\int_{L} \frac{-y d x+x d y}{x^{2}+y^{2}}$ 的值为 \xzanswer{A} 


\fourchoices
{恒为 $-\pi$}
{恒为 $ 0 $}
{恒为 $\pi$}
{与曲线 $L$ 有关}

\item 
广义积分 $\int_{0}^{+\infty} \frac{\sin x^{2}}{x^{p}} d x$ 的收敛域为 \xzanswer{D} 


\fourchoices
{$p>-1$}
{$0<p<3$}
{$-1<p<1$}
{$-1<p<3$}


	
	
\end{enumerate}

\banswer{
\begin{enumerate}
	%\renewcommand{\labelenumi}{\arabic{enumi}.}
	\item
	\textbf{C}。因 $\arctan\frac1x$ 有界，$x\to0$ 时
	$|x|^{a}\arctan\frac1x\to0=f(0)$ 当且仅当 $a>0$，这是连续的充要条件。
	再看差商：$x\to0^{+}$ 时
	\[
	\frac{f(x)-f(0)}{x}=x^{a-1}\arctan\frac1x\sim\frac\pi2x^{a-1} ,
	\]
	$x\to0^{-}$ 时 $\dfrac{|x|^{a}\arctan(1/x)}{x}=-|x|^{a-1}\arctan\frac1x
	\sim\frac\pi2|x|^{a-1}$，两侧同型。于是 $a>1$ 时导数存在且为 $0$；
	$a=1$ 时左右差商都趋于 $\frac\pi2$，仍可导；只有 $0<a<1$ 时差商趋于
	$+\infty$，连续而不可导。

	\item
	原卷此处题面缺失（仅存题号），无从作答。

	\item
	\textbf{C}。记部分和 $S_{n}=\sum\limits_{k=1}^{n}a_{k}$，所述条件即
	\[
	\forall\varepsilon>0,\ \exists N,\ \forall n>N:\quad
	\left|S_{n}-S_{N}\right|<\varepsilon .
	\]
	若它成立，则对一切 $m,n>N$ 有
	$\left|S_{n}-S_{m}\right|\leqslant\left|S_{n}-S_{N}\right|
	+\left|S_{m}-S_{N}\right|<2\varepsilon$，
	即 $\{S_{n}\}$ 为 Cauchy 列，级数收敛，故条件充分；
	反之级数收敛时，由 Cauchy 收敛准则该条件自然成立，故亦必要。
	所以是充分必要条件。

	\par\vspace{1mm}\noindent
	\textbf{注：}{\kaishu
	原题面末项印作 $a_{N}$，与前面的 $a_{N+1},a_{N+2},\cdots$ 矛盾，
	应为 $a_{n}$，已改正。
	}\par

	\item
	\textbf{B}。取 $x=R\cos\theta$、$y=R\sin\theta$、$0\leqslant z\leqslant R$，
	则 $dS=R\,d\theta\,dz$，于是
	\begin{align*}
	\iint_{S}\left(x^{2}+z^{2}\right)dS
	&=R\int_{0}^{2\pi}\!\!\int_{0}^{R}
	\left(R^{2}\cos^{2}\theta+z^{2}\right)dz\,d\theta\\
	&=R\left(\pi R^{3}+2\pi\cdot\frac{R^{3}}{3}\right)
	=\frac53\pi R^{4} .
	\end{align*}

	\item
	\textbf{A}。在不含原点的单连通区域内
	\[
	\frac{-y\,dx+x\,dy}{x^{2}+y^{2}}=d\left(\arctan\frac yx\right)=d\theta
	\]
	是全微分（极角的微分）。$L$ 除端点外全在上半平面，
	故沿 $L$ 极角由 $\theta=\pi$ 连续变到 $\theta=0$，积分值
	$=\Delta\theta=0-\pi=-\pi$，与 $L$ 的具体形状无关。

	\item
	\textbf{D}。令 $t=x^{2}$，则 $dx=\dfrac{dt}{2\sqrt t}$，
	\[
	\int_{0}^{+\infty}\frac{\sin x^{2}}{x^{p}}\,dx
	=\frac12\int_{0}^{+\infty}t^{-\frac{p+1}{2}}\sin t\,dt .
	\]
	$t\to0^{+}$ 时被积函数 $\sim\frac12t^{\frac{1-p}{2}}$，收敛要求
	$\frac{1-p}{2}>-1$，即 $p<3$；$t\to+\infty$ 时由 Dirichlet 判别法，
	需 $t^{-\frac{p+1}{2}}$ 单调趋于 $0$，即 $p>-1$
	（$p=-1$ 时 $\int^{+\infty}\sin t\,dt$ 发散）。故收敛域为 $-1<p<3$。
\end{enumerate}
}

\item
(本题共 5 小题，每小题 8 分，满分 40 分)
\begin{enumerate}
	%\renewcommand{\labelenumi}{\arabic{enumi}.}
	% A(\Alph) a(\alph) I(\Roman) i(\roman) 1(\arabic)
	%设定全局标号series=example	%引用全局变量resume=example
	%[topsep=-0.3em,parsep=-0.3em,itemsep=-0.3em,partopsep=-0.3em]
	%可使用leftmargin调整列表环境左边的空白长度 [leftmargin=0em]
	\item
求 $\int \max (x, 1) d x$。
	
	
	\item
	计算无穷积分 $\int_{1}^{+\infty} \frac{1}{x\left(x^{m}+1\right)} d x$, 其中 $m$ 为正整数。
	
	
	\item
	设 $G_{n}=\sqrt[n]{(n+1)(n+2) \cdots(n+n)}$, 求 $\lim\limits _{n \rightarrow \infty} \frac{G_{n}}{n}$ 。
	
	
	\item
	证明：当 $x>0$ 的时候，有 $\left(1+\frac{1}{x}\right)^{x}<e<\left(1+\frac{1}{x}\right)^{x+1}$ 。
	
	
	\item
	函数 $f(x)$ 在 $[0,+\infty)$ 上有一阶连续导函数, 对所有 $x \geqslant 0$, 有 $f(x) \leqslant e^{-x}$, 且 $f(0)=1$ 。 证明：存在 $\xi>0$, 使得 $f^{\prime}(\xi)=-e^{-\xi}$。
	
	
	
	
\end{enumerate}

	
\banswer{
	\begin{enumerate}
		%\renewcommand{\labelenumi}{\arabic{enumi}.}
		% A(\Alph) a(\alph) I(\Roman) i(\roman) 1(\arabic)
		%设定全局标号series=example	%引用全局变量resume=example
		%[topsep=-0.3em,parsep=-0.3em,itemsep=-0.3em,partopsep=-0.3em]
		%可使用leftmargin调整列表环境左边的空白长度 [leftmargin=0em]
		\item
		$\int \max (x, 1) d x= 
		\left\{ 
\begin{array}{lr}
	x +C & x \leq 1\\
	\frac{x^{2}+1}{2} +C &x \geq 1
\end{array}\right.
$
\item 		
$ \frac{1}{m} \ln 2 $
\item 
$ \frac{4}{e} $	
\item 
\textbf{证} 当 $x>0$ 时，$\displaystyle\ln\!\left(1+\frac{1}{x}\right)=\ln\frac{x+1}{x}=\int_{x}^{x+1}\frac{dt}{t}$。由 $\frac{1}{t}$ 在 $[x,x+1]$ 上严格递减得
\[
\frac{1}{x+1}<\int_{x}^{x+1}\frac{dt}{t}<\frac{1}{x}.
\]
即 $\frac{1}{x+1}<\ln\!\left(1+\frac{1}{x}\right)<\frac{1}{x}$。乘以 $x$ 得 $\frac{x}{x+1}<x\ln\!\left(1+\frac{1}{x}\right)<1$；乘以 $x+1$ 得 $1<(x+1)\ln\!\left(1+\frac{1}{x}\right)$。取指数即得
\[
\left(1+\frac{1}{x}\right)^{x}<e<\left(1+\frac{1}{x}\right)^{x+1}.
\]
\item 
\textbf{解} 本题结论不成立。反例：取 $f(x)=e^{-x}-x\;(x\geqslant0)$，则 $f$ 在 $[0,+\infty)$ 上有一阶连续导数，$f(0)=1$，且 $f(x)\leqslant e^{-x}$。但
\[
f'(x)=-e^{-x}-1,\qquad -e^{-\xi}=-e^{-\xi},
\]
方程 $f'(\xi)=-e^{-\xi}$ 即 $-1=0$，无解。故不存在 $\xi>0$ 满足要求。
	\end{enumerate}
	
	
}
	
	
\item 
(本题共 3 小题，每小题 12 分，满分 36 分)
\begin{enumerate}
	%\renewcommand{\labelenumi}{\arabic{enumi}.}
	% A(\Alph) a(\alph) I(\Roman) i(\roman) 1(\arabic)
	%设定全局标号series=example	%引用全局变量resume=example
	%[topsep=-0.3em,parsep=-0.3em,itemsep=-0.3em,partopsep=-0.3em]
	%可使用leftmargin调整列表环境左边的空白长度 [leftmargin=0em]
	\item
 求函数 $z=x^{2} y(3-x-y)$ 在闭区域 $D: x \geqslant 0, y \geqslant 0, x+y \leqslant 4$ 上的最大值和最小值。
	
	
	\item
	设 $f(x)=\cos x+\frac{1}{4} \int_{0}^{2 \pi}(2 x-t) f\left(\frac{t}{2}\right) d t$, 其中 $f(x)$ 为连续函数，求 $f(x)$ 。
	
	
	\item
	设 $a, b, c>0$, 求曲面 $x^{2}+y^{2}+\frac{a^{2}-b^{2}}{c^{2}} z^{2}=a^{2}$ 与 $|z| \leqslant c$ 所截物体的体积。
	
	
	
	
\end{enumerate}



\banswer{
	\begin{enumerate}
		%\renewcommand{\labelenumi}{\arabic{enumi}.}
		% A(\Alph) a(\alph) I(\Roman) i(\roman) 1(\arabic)
		%设定全局标号series=example	%引用全局变量resume=example
		%[topsep=-0.3em,parsep=-0.3em,itemsep=-0.3em,partopsep=-0.3em]
		%可使用leftmargin调整列表环境左边的空白长度 [leftmargin=0em]
		\item
		$ z_{\max}=\frac{81}{64}  \quad ( \frac{ 3 }{ 2 } , \frac{ 3 }{ 4 } )  $， \quad $ z_{\min}=-\frac{256}{27}  \quad ( \frac{ 8 }{ 3 } , \frac{ 4 }{ 3 } ) $
		\item
		令 $u=\frac t2$，则
		\[
		\frac14\int_{0}^{2\pi}(2x-t)f\!\left(\frac t2\right)dt
		=\int_{0}^{\pi}(x-u)f(u)\,du ,
		\]
		故 $f(x)=\cos x+Px-Q$，其中
		$P=\int_{0}^{\pi}f(u)du$、$Q=\int_{0}^{\pi}uf(u)du$ 均为常数。
		把这一形式代回（用 $\int_{0}^{\pi}\cos u\,du=0$、
		$\int_{0}^{\pi}u\cos u\,du=-2$）得
		\[
		P=\frac{\pi^{2}}{2}P-\pi Q,\qquad
		Q=-2+\frac{\pi^{3}}{3}P-\frac{\pi^{2}}{2}Q ,
		\]
		解出 $P=\dfrac{24\pi}{\pi^{4}+12}$、
		$Q=\dfrac{12\left(\pi^{2}-2\right)}{\pi^{4}+12}$，所以
		\[
		f(x)=\cos x+\frac{24\pi}{\pi^{4}+12}x
		-\frac{12\left(\pi^{2}-2\right)}{\pi^{4}+12} .
		\]

		\par\vspace{1mm}\noindent
		\textbf{注：}{\kaishu
		原答案作 $f(x)=\frac12\cos x+Ae^{x}+Be^{-x}$，有误。
		本题积分限 $0$ 到 $2\pi$ 是常数限，$f$ 由等式唯一确定，不含任意常数；
		带 $Ae^{x}+Be^{-x}$ 的结果属于变上限方程
		$f(x)=\cos x+\int_{0}^{x}(x-t)f(t)dt$（求导两次给出 $f''-f=-\cos x$）。
		上面的 $f$ 已用数值积分逐点校验，等式两端完全相符。
		}\par
		\item
		$ V=\frac{2\pi c}{3}(2a^{2}+b^{2}) $
	\end{enumerate}
	
	
}


\newpage
\item 
(本题共 2 小题，每小题 12 分，满分 24 分)
\begin{enumerate}
	%\renewcommand{\labelenumi}{\arabic{enumi}.}
	% A(\Alph) a(\alph) I(\Roman) i(\roman) 1(\arabic)
	%设定全局标号series=example	%引用全局变量resume=example
	%[topsep=-0.3em,parsep=-0.3em,itemsep=-0.3em,partopsep=-0.3em]
	%可使用leftmargin调整列表环境左边的空白长度 [leftmargin=0em]
	\item
	设 $f(x)=\left\{\begin{aligned}&\frac{\pi-1}{2} x \ & \ 0 \leqslant x \leqslant 1 \\ &\frac{\pi-x}{2} \ & \ 1<x \leqslant \pi\end{aligned}\right.$，将 $f(x)$ 展开为周期为 $2 \pi$ 的正弦级数, 并求$\sum\limits_{n=1}^{\infty} \frac{\sin ^{2} n}{n^{2}}$ 和 $\sum\limits_{n=1}^{\infty} \frac{\sin ^{2} n}{n^{4}}$。
	
	\item 
设区域 $c$ 由曲面 $2 x^{2}+y^{2}-z^{2}=1$ 及平面 $z=1, z=-1$ 所围成, $S$ 为 $c$ 的全表 面外侧， 又设 $\vec{v}=\left(2 x^{2}+y^{2}+z^{2}\right)^{-\frac{3}{2}}(x \vec{i}+y \vec{j}+z \vec{k})$ 。
\begin{enumerate}
	%\renewcommand{\labelenumi}{\arabic{enumi}.}
	% A(\Alph) a(\alph) I(\Roman) i(\roman) 1(\arabic)
	%设定全局标号series=example	%引用全局变量resume=example
	%[topsep=-0.3em,parsep=-0.3em,itemsep=-0.3em,partopsep=-0.3em]
	%可使用leftmargin调整列表环境左边的空白长度 [leftmargin=0em]
	\item
求 $\operatorname{div} \vec{v}$。
	\item 
	求积分 $\iint_{S} \frac{x d y d z+y d z d x+z d x d y}{\left(2 x^{2}+y^{2}+z^{2}\right)^{\frac{3}{2}}}$。
\end{enumerate}
	
	
	
	
\end{enumerate}


\banswer{
	\begin{enumerate}
		%\renewcommand{\labelenumi}{\arabic{enumi}.}
		% A(\Alph) a(\alph) I(\Roman) i(\roman) 1(\arabic)
		%设定全局标号series=example	%引用全局变量resume=example
		%[topsep=-0.3em,parsep=-0.3em,itemsep=-0.3em,partopsep=-0.3em]
		%可使用leftmargin调整列表环境左边的空白长度 [leftmargin=0em]
		\item
		$ f(x)=\sum\limits_{n=1}^{\infty} \frac{\sin n}{n^{2}} \sin nx (0 \leq x \leq \pi)$，利用$ x=1 $得到$\sum\limits_{n=1}^{\infty} \frac{\sin ^{2} n}{n^{2}}=\frac{\pi -1}{2}$，再由帕塞瓦尔恒等式(Parseval's identity)，得到 $\sum\limits_{n=1}^{\infty} \frac{\sin ^{2} n}{n^{4}}=\frac{1}{\pi} \int_{-\pi}^{\pi} f^{2}(x) dx = \frac{(\pi-1)^{2}}{6} $
		\item 
		\begin{enumerate}
			%\renewcommand{\labelenumi}{\arabic{enumi}.}
			% A(\Alph) a(\alph) I(\Roman) i(\roman) 1(\arabic)
			%设定全局标号series=example	%引用全局变量resume=example
			%[topsep=-0.3em,parsep=-0.3em,itemsep=-0.3em,partopsep=-0.3em]
			%可使用leftmargin调整列表环境左边的空白长度 [leftmargin=0em]
			\item
			$\operatorname{div} \vec{v}=0$
			\item 
			$ 2\sqrt{2} \pi $
			
		\end{enumerate}
		
		
		
	\end{enumerate}
	
	
}
	
	
\end{enumerate}


